This thesis presents the design and implementation of a business intelligence prototype that integrates credit data processing, multi-horizon prediction, experiment tracking, and analytical visualization for an Ecuadorian social and solidarity economy institution. PostgreSQL is used for storage and the data mart, Apache Airflow orchestrates training and inference workflows, MLflow records experiments, and Apache Superset provides interactive dashboards. The evaluation compares LightGBM, a convolutional neural network, and a multilayer perceptron across eighteen forecasting horizons. In the reported results, LightGBM achieved a mean ROC-AUC of 0.80, whereas the CNN and MLP implementations yielded a ROC-AUC of 0.50. LightGBM performance decreased as the horizon increased, and recall fell to 0.034 at eighteen months, which limits the operational interpretation of long-term predictions. Under the evaluated configuration, LightGBM showed the highest discriminative performance; however, the lack of repeated runs, statistical tests, external validation, and complete documentation of the heuristic target prevents claims of general superiority or causality. The main contribution is the technical integration of data processing, modeling, experiment tracking, analytical storage, and visualization into a functional prototype that requires further methodological and operational validation.
\\
\\
\noindent\textbf{Keywords}:
credit risk; multi-horizon prediction; LightGBM; neural networks; business intelligence; MLOps; data mart; Apache Airflow.